\documentclass[12pt]{article}
\usepackage[utf8]{inputenc}
\usepackage[T1]{fontenc}
\usepackage{hyperref}
\title{Onya Darat Language Resources for K-12 Curriculum Development}
\author{Compiled by Claude Sonnet 5.5 (claude-sonnet-5-5) for Andr\'as Kornai}
\date{}
\begin{document}
\maketitle

\section{Introduction}
Onya Darat (ISO 639-3 \texttt{sdm}; Glottocode sema1269), also called Semandang, is a Land Dayak (Bidayuhic) language of Ketapang regency, West Kalimantan, Indonesia. Wikipedia gives about 20,000 speakers (2004). It is among the least-studied Austronesian languages, and no school materials were found.

\section{Primary Educational Resources}
No primer, textbook or school programme for Onya Darat turned up in the searches.

\section{Linguistic and Reference Materials}
\begin{itemize}
\item Wikipedia, Semandang language (Onya Darat), with dialects Kualan, Samandang, Baram and Simpang. \url{https://en.wikipedia.org/wiki/Semandang_language}
\item Max Planck EVA project ``Documentation of the phonetic structures of Onya Darat''. \url{https://www.eva.mpg.de/linguistics/past-research-resources/documentation-and-description/documentation-of-the-phonetic-structures-of-onya-darat/}
\item Wacana (Universitas Indonesia) journal item found in search; reported by the search summary to relate to the Onya Darat orthography article by Tadmor (2009). Its exact content was not verified. \url{https://scholarhub.ui.ac.id/wacana/vol11/iss2/5}
\item Glottolog entry. \url{https://glottolog.org/resource/languoid/id/sema1269}
\end{itemize}

\section{Religious and Community Resources}
A search for a Bible translation found nothing for this language.

\section{Digital and Audio Resources}
The EVA project page above documents phonetic recordings work; the searches did not confirm public downloadable audio.

\section{Curriculum Development Context}
Only a few academic descriptions exist, and no community literacy materials were found. The Tadmor (2009) orthography article, if it is the Wacana item, would be the starting point for a writing system. Any curriculum would need to be built with Onya Darat speakers in Ketapang.

\end{document}